Spectre attacks exploit hardware mispredictions to coerce victim programs into leaking secrets over hardware side channels via transient (i.e., uncommitted) instructions. To provide software with a specific security guarantee against these attacks at minimal performance cost, many Spectre defenses require changes to conventional processor microarchitectures. These changes are usually prototyped in microarchitectural simulators like gem5, which can evaluate a defense implementation's performance and functional correctness by running sizable programs, but offer no analogous mechanism for checking whether it upholds its security guarantee. As a result, subtle security bugs can go undetected, voiding a defense's security guarantee and potentially understating its performance overhead.

We propose \toolname{}, a three-part framework for automated, thorough security validation of simulated Spectre defenses. First, a new specification language enables formalizing a Spectre defense's security guarantee as a predicate identifying which (leaky) program executions the defense forbids. Second, a synthesizer uses such a specification to generate a compact, targeted suite of \textit{security litmus tests}---small programs, each paired with a target execution (i.e., specific microarchitectural control-/data-flows) each exhibiting forbidden leakage. Third, an executor actively steers each test toward its target execution on a simulated defense prototype (here, in gem5) and checks whether the defense prevents it. Applied to gem5 prototypes of five publicly available Spectre defenses and one obtained from its authors, \toolname{} uncovers all bugs that were found by prior state-of-the-art simulator tester AMuLeT, and are in-scope for the defense, plus six new bugs. Compared to AMuLeT, SimSpect needs at least $23\times$ fewer tests to find each known bug.

% We propose \textbf{SimSpect}, a rigorous, targeted, and automated approach for validating early-stage simulator prototypes of hardware support for side-channel defenses. We present a formal language intermediate representation for specifying leakage prohibited by a defense as well as a validator that generates and executes a comprehensive litmus test suite from this specification, combining into a pipeline for rigorous defense validation and iteration. 
% Applying SimSpect to all accessible simulator prototypes of state-of-the-art defenses, we replicate all in-scope previously known security bugs, and find the following previously unknown bugs: 
% (1) STT mistakenly fails to taint one register, 
% ReCon's underlying implementation of STT (2) skips a taint check on store-to-load forwarding and (3) tracks only the oldest taint source as opposed to all taint sources, 
% SpecLFB (4) misimplements its ROB walk order and (5) fails to protect under the shadow of a jmp instruction, and finally 
% (6) there is a visibility-point/squash timing race shared by STT and SPT.
% Compared to prior work, SimSpect decreases the number of tests needed to find known bugs on average by a factor of 6.5\anna{update}.


% find a bug in STT where it fails to taint an FPR, uncover a previously unknown visibility-point/squash timing race shared by STT and SPT, and find two bugs in ReCon's underlying implementation of STT: store-to-load forwarding bypasses its taint check, and an instruction that merges two tainted operands keeps only the older operand's taint. Finally, we show that our test generation corpus takes less time to rediscover known bugs, for instance, on SpecLFB, SimSpect reaches the known first-load bug in 6.5x fewer tests than AMuLeT.



% Hardware side-channel attacks pass secrets to instructions with data-dependent optimizations that leak victim data via timing or other indicators. Proposed defenses against these attacks include hardware modifications that are prototyped in computer architecture simulators such as gem5. However, these prototyped mechanisms frequently contain security bugs that invalidate their reported performance results, and there exists no comprehensive early-stage validation method for their proposed security guarantees. 
% Many hardware defenses have been proposed to mitigate the dangers posed by side-channel attacks, particularly Spectre attacks but they present performance guarantees for their hardware via simulator prototypes. There exists no comprehensive validation method for the security guarantees of such simulators, and most have been shown to contain security-violating bugs that invalidate performance results. We develop SimSpect, a rigorous, targeted, and automated approach for validating early-stage simulator prototypes of hardware support for side-channel defenses (SCDs). We present a formal language intermediate representation for specifying hardware SCDs and proving security properties as well as a validator that generates and executes a comprehensive litmus test suite from this specification, together forming a pipeline for rigorous defense validation and iteration.


% security bugs \anna{that violate their defense's proposed security guarantees} as well as \textit{six new ones}.
% \urgentbr
% \urgentnote{``All known security bugs'' needs its scope and count. AMuLeT reports nine bugs on these defenses (fu-2025-amulet.pdf p1); five are inside each defense's own guarantee (UV1, UV3, UV4, UV5, UV6) and SimSpect finds all five (UV3 only with the read-modify-write corpus and a cold cache, 407 of 590 tests vs 0 of 590 warm; UV5 in one test so far). The other four are out of scope by each paper (UV2, KV1, KV2, KV3). On ``six new ones'': you decided SpecLFB's walk order is UV6 reaching further than AMuLeT published, yet the abstract and the intro count it as new bug (4). Options: count five new bugs plus a wider form of UV6, or keep six and call (4) a new form of UV6. \textbf{PROPOSED TEXT} (edits the sentence just before this note; recommendation, using the first option): ``\ldots{} \toolname{} uncovers all five previously known bugs that violate their defense's own guarantee, as well as five new bugs and a wider form of a known one.'' \textit{Source:} simspect wiki/findings/mdsToRead/AMULET\_BUGS\_DETECTION\_ACCOUNTING.md sec.~2; URGENT\_TODO.md walk-order decision (2026-11-09).}
% \notebr